% problem-000421 5 脂肪酸钠缔合胶体
\section*{5. 脂肪酸钠缔合胶体}
\textit{下列为分问。}

\subsection*{5-1}
根据上述事实，你认为McBain的推论是否合理？说明理由。

\subsection*{5-2}
McBain对脂肪酸钠体系进⾏了量热实验，发现在拐点浓度附近，由拐点前到拐点后体系的标准焓变⼏乎为零，这让他百思不得其解。经过⻓期的研究和思考，他⿎起勇⽓在国际学术会议上宣讲了⾃⼰的研究结果，

认为脂肪酸在拐点浓度后的缔合胶体形成是热⼒学稳定的。当时的会议主席对此勃然⼤怒，没有等报告讲完就以“McBain，胡说！”，将McBain轰下讲台。

假设在脂肪酸盐拐点浓度前、后体系中⽔的结构性质没有变化，结合基础热⼒学和熵的统计意义，简单说明会议主席认为McBain的观点“荒谬”的原因。

\subsection*{5-3}
然⽽，⼤量研究结果证实了McBain的结论是正确的，原因是“在脂肪酸盐拐点浓度前、后⽔的结构性质没有变化”，这⼀假设是错误。如果认为“脂肪酸盐在拐点浓度前、后对溶剂⽔的结构会产⽣不同的影响”，那么请基于熵的统计意义，推测⽔的结构在拐点浓度前和拐点后哪个更有序，说明理由。

\subsection*{5-4}
除脂肪酸盐外还有很多物质具有类似的变化规律。这些物质（表⾯活性剂）的分⼦结构具有⼀个共同的特点，都有亲⽔头基和蔬⽔尾链（如下图所⽰）。请给出表⾯活性剂在⽔中形成的“缔合胶体”的结构⽰意图。若表⾯活性剂的分⼦⻓度为1.0nm，估算“缔合胶体”在刚过拐点浓度时的最⼤尺⼨。

（图：）
